\subsection{Haces coherentes}
\label{subsection:0.5.3-fr}

\begin{env}[5.3.1]
\label{0.5.3.1-fr}
Se dice que un $\mathcal{O}_X$-Módulo $\mathcal{F}$ es \emph{coherente} si
verifica las dos condiciones siguientes:
\begin{enumerate}
  \item[$a)$] $\mathcal{F}$ es de tipo finito.
  \item[$b)$] Para todo abierto $U\subset X$, todo entero $n>0$, y todo
    homomorfismo
    $u:\mathcal{O}_X^n|U\to\mathcal{F}|U$, el núcleo de $u$ es de tipo finito.
\end{enumerate}
Se observará que estas dos condiciones son de carácter \emph{local}.

Para la mayoría de las demostraciones de las propiedades de los haces coherentes
que se recuerdan a continuación, cf. (FAC), I, 2.
\end{env}

\begin{env}[5.3.2]
\label{0.5.3.2-fr}
Todo $\mathcal{O}_X$-Módulo coherente admite una presentación finita (5.2.5); el
recíproco no es necesariamente exacto, pues el propio $\mathcal{O}_X$ no es
necesariamente un $\mathcal{O}_X$-Módulo coherente.

Todo sub-$\mathcal{O}_X$-Módulo \emph{de tipo finito} de un
$\mathcal{O}_X$-Módulo coherente es coherente; toda suma directa \emph{finita}
de $\mathcal{O}_X$-Módulos coherentes es un $\mathcal{O}_X$-Módulo coherente.
\end{env}

\begin{env}[5.3.3]
\label{0.5.3.3-fr}
Si $0\to\mathcal{F}\to\mathcal{G}\to\mathcal{H}\to0$ es una sucesión exacta de
$\mathcal{O}_X$-Módulos y si dos de estos $\mathcal{O}_X$-Módulos son
coherentes, lo mismo ocurre con el tercero.
\end{env}

\begin{env}[5.3.4]
\label{0.5.3.4-fr}
Si $\mathcal{F}$ y $\mathcal{G}$ son dos $\mathcal{O}_X$-Módulos coherentes,
$u:\mathcal{F}\to\mathcal{G}$ un homomorfismo, $\operatorname{Im}(u)$,
$\operatorname{Ker}(u)$ y $\operatorname{Coker}(u)$ son
$\mathcal{O}_X$-Módulos coherentes. En particular, si $\mathcal{F}$ y
$\mathcal{G}$ son sub-$\mathcal{O}_X$-Módulos coherentes de un
$\mathcal{O}_X$-Módulo coherente, $\mathcal{F}+\mathcal{G}$ y
$\mathcal{F}\cap\mathcal{G}$ son coherentes.
\end{env}

\begin{env}[5.3.5]
\label{0.5.3.5-fr}
Si $\mathcal{F}$ y $\mathcal{G}$ son dos $\mathcal{O}_X$-Módulos coherentes,
lo mismo ocurre con $\mathcal{F}\otimes_{\mathcal{O}_X}\mathcal{G}$ y con
$\mathcal{H}\!om_{\mathcal{O}_X}(\mathcal{F},\mathcal{G})$.
\end{env}

\begin{env}[5.3.6]
\label{0.5.3.6-fr}
Sean $\mathcal{F}$ un $\mathcal{O}_X$-Módulo coherente, $\mathcal{J}$ un
haz coherente de ideales de $\mathcal{O}_X$. Entonces el
$\mathcal{O}_X$-Módulo $\mathcal{J}\mathcal{F}$ es coherente, como imagen
de $\mathcal{J}\otimes_{\mathcal{O}_X}\mathcal{F}$ por el homomorfismo
canónico
$\mathcal{J}\otimes_{\mathcal{O}_X}\mathcal{F}\to\mathcal{F}$
(5.3.4 y 5.3.5).
\end{env}

\begin{env}[5.3.7]
\label{0.5.3.7-fr}
Se dice que una $\mathcal{O}_X$-Álgebra $\mathcal{A}$ es \emph{coherente} si
$\mathcal{A}$ es un $\mathcal{O}_X$-Módulo coherente. En particular,
$\mathcal{O}_X$ es un \emph{haz coherente de anillos} si y solo si,
para todo abierto $U\subset X$ y todo homomorfismo de la forma
$u:\mathcal{O}_X^p|U\to\mathcal{O}_X|U$, el núcleo de $u$ es un
$(\mathcal{O}_X|U)$-Módulo de tipo finito.

Si $\mathcal{O}_X$ es un haz coherente de anillos, todo
$\mathcal{O}_X$-Módulo $\mathcal{F}$ que admite una presentación finita (5.2.5)
es coherente, en virtud de (5.3.4).

El \emph{anulador} de un $\mathcal{O}_X$-Módulo $\mathcal{F}$ es el núcleo
$\mathcal{J}$ del homomorfismo canónico
$\mathcal{O}_X\to\mathcal{H}\!om_{\mathcal{O}_X}(\mathcal{F},\mathcal{F})$
que a toda sección $s\in\Gamma(U,\mathcal{O}_X)$ hace corresponder la
multiplicación por $s$ en
$\operatorname{Hom}(\mathcal{F}|U,\mathcal{F}|U)$; si $\mathcal{O}_X$ es
coherente y si $\mathcal{F}$ es un $\mathcal{O}_X$-Módulo coherente,
$\mathcal{J}$ es coherente (5.3.4 y 5.3.5) y para todo $x\in X$,
$\mathcal{J}_x$ es el anulador de $\mathcal{F}_x$ (5.2.6).
\end{env}

\oldpage[0\textsubscript{I}]{48}

\begin{env}[5.3.8]
\label{0.5.3.8-fr}
Supongamos que $\mathcal{O}_X$ es coherente; sean $\mathcal{F}$ un
$\mathcal{O}_X$-Módulo coherente, $x$ un punto de $X$, $M$ un submódulo de
tipo finito de $\mathcal{F}_x$; existe entonces una vecindad abierta $U$ de $x$
y un sub-$(\mathcal{O}_X|U)$-Módulo coherente $\mathcal{G}$ de
$\mathcal{F}|U$ tal que $\mathcal{G}_x=M$ (T, 4.1, lema 1).

Este resultado, unido a las propiedades de los sub-$\mathcal{O}_X$-Módulos de un
$\mathcal{O}_X$-Módulo coherente, impone condiciones necesarias a los anillos
$\mathcal{O}_x$ para que $\mathcal{O}_X$ sea coherente. Por ejemplo (5.3.4),
la intersección de dos ideales de tipo finito de $\mathcal{O}_x$ debe ser también
un ideal de tipo finito.
\end{env}

\begin{env}[5.3.9]
\label{0.5.3.9-fr}
Supongamos que $\mathcal{O}_X$ es coherente, y sea $M$ un $\mathcal{O}_x$-módulo
que admite una presentación finita, por tanto isomorfo a un conúcleo de un
homomorfismo $\varphi:\mathcal{O}_x^p\to\mathcal{O}_x^q$; entonces existe una
vecindad abierta $U$ de $x$ y un $(\mathcal{O}_X|U)$-Módulo coherente
$\mathcal{F}$ tal que $\mathcal{F}_x$ es isomorfo a $M$. En efecto, en virtud
de (5.2.6), existe una sección $u$ de
$\mathcal{H}\!om_{\mathcal{O}_X}(\mathcal{O}_X^p,\mathcal{O}_X^q)$ tal que $u_x=\varphi$; el conúcleo
$\mathcal{F}$ del homomorfismo
$u:\mathcal{O}_X^p|U\to\mathcal{O}_X^q|U$ responde a la cuestión (5.3.4).
\end{env}

\begin{env}[5.3.10]
\label{0.5.3.10-fr}
Supongamos que $\mathcal{O}_X$ es coherente, y sea $\mathcal{J}$ un haz coherente
de ideales de $\mathcal{O}_X$. Para que un
$(\mathcal{O}_X/\mathcal{J})$-Módulo $\mathcal{F}$ sea coherente, es necesario y
suficiente que sea coherente como $\mathcal{O}_X$-Módulo. En particular,
$\mathcal{O}_X/\mathcal{J}$ es un haz coherente de anillos.
\end{env}

\begin{env}[5.3.11]
\label{0.5.3.11-fr}
Sea $f:X\to Y$ un morfismo de espacios anillados, y supongamos que
$\mathcal{O}_X$ es coherente; entonces, para todo $\mathcal{O}_Y$-Módulo coherente
$\mathcal{G}$, $f^*(\mathcal{G})$ es un $\mathcal{O}_X$-Módulo coherente. En
efecto, con las notaciones de (5.2.4), se puede suponer que $\mathcal{G}|V$ es
conúcleo de un homomorfismo
$v:\mathcal{O}_Y^q|V\to\mathcal{O}_Y^p|V$; como $f_U^*$ es exacto por la derecha,
$f^*(\mathcal{G})|U=f_U^*(\mathcal{G}|V)$ es conúcleo del homomorfismo
$f_U^*(v):\mathcal{O}_X^q|U\to\mathcal{O}_X^p|U$, de donde nuestra afirmación.
\end{env}

\begin{env}[5.3.12]
\label{0.5.3.12-fr}
Sean $Y$ una parte cerrada de $X$, $j:Y\to X$ la inyección canónica,
$\mathcal{O}_Y$ un haz de anillos sobre $Y$, y pongamos
$\mathcal{O}_X=j_*(\mathcal{O}_Y)$. Para que un $\mathcal{O}_Y$-Módulo
$\mathcal{G}$ sea de tipo finito (resp. cuasicoherente, coherente), es necesario y
suficiente que $j_*(\mathcal{G})$ sea un $\mathcal{O}_X$-Módulo de tipo finito
(resp. cuasicoherente, coherente).
\end{env}

